\subsection{G 构建 SWE 环境的基础设施}\label{g-infrastructure-for-building-swe-environments}

该基础设施的设计主要旨在应对规模、稳定性和安全方面的挑战。构建基础设施组织为一个双池（two-pool）Ray 集群，运行在约 30,000 个 CPU 核心上。主池（main pool）负责工作调度、Lance I/O 和流水线状态跟踪。构建池（builder pool）运行实际的容器构建；所有构建节点都标记了 Ray 自定义资源标志（custom resource flag），构建步骤利用 Ray 的逻辑资源调度，确保其 actor 只落在具备构建能力的节点上。

两个池之所以分离，是因为容器构建对资源要求高且容易失败（磁盘耗尽、OOM、构建挂起），而协调工作轻量且无状态；将它们隔离可确保构建 pod 的失败不会干扰编排和状态管理。这种分离还实现了凭据隔离：构建 pod 执行来自开源仓库的不可信 Dockerfile，不持有流水线数据存储或内部服务的凭据。Pod 失败——无论由磁盘压力、硬件故障还是节点驱逐引起——都由 Kubernetes 重启处理；重启后的 pod 重新连接到 Ray 集群，并在下一个主池 worker 调度新任务时被自动接管。

每个构建 pod 包含两个容器，共享基于 NVMe 的本地存储，协同部署以提升构建与评分的效率。主容器运行 Ray worker，并使用 rootless podman 加载构建好的镜像、运行评分容器。sidecar 容器运行 rootless BuildKit 守护进程，通过 buildctl 进行 OCI 镜像构建。两个容器通过共享的 NVMe 卷通信——BuildKit 将 OCI tar 归档写入该卷，podman 再加载这些归档——并通过 localhost 走 BuildKit 控制通道。这种协同部署避免了大型容器镜像在构建与评分之间的网络传输；整个构建-加载-运行循环都在本地 NVMe 上完成。Init 容器负责一次性初始化，包括通过联合身份（federated identity）完成容器镜像仓库的身份验证。

该流水线还利用 BuildKit 分层缓存，减少同一仓库内不同问题之间的冗余工作。问题按仓库分组，并分派给构建节点上同一个持久的 Ray actor。由于该 actor 在同一仓库的问题之间保持存活，共享的基础层——apt-get install、pip install、依赖编译——只构建一次并被复用，因为依赖安装往往主导 Docker 构建



表 15. 来自我们分类体系的高层约束示例，按客观（硬性）和主观（软性）类别划分，并给出相应的子类别与示例。

{\def\LTcaptype{none} % do not increment counter
\begin{longtable}[]{@{}|l|l|@{}}
\toprule\noalign{}
\endhead
\bottomrule\noalign{}
\endlastfoot
\hline
约束类型 & 子类别 \& 示例 \\
\hline
客观（硬性）约束 & \\
\hline
数值 & 显式长度限制、句子/段落数量、条目数量、范围约束和比例。 \\
\hline
格式 & 文档结构、列表、表格、代码格式、数据序列化（JSON/YAML）、
标记/样式、模板和引用格式。 \\
\hline
内容 & 必需/禁止的术语和符号、必需的元素或概念、来源/引用要求，
以及时间/地域范围。 \\
\hline
语言 & 语言选择、方言、语域、句子结构、词性、形态约束和语法规则。 \\
\hline
时间 & 顺序 & 时间先后排序、程序步骤和字母顺序。 \\
\hline
否定约束 & 排除特定行为或格式模式的指令 \\
\hline
主观（软性）约束 & \\
\hline
风格 & 语气、口吻、词汇水平、受众适配、情感效价和品牌语调。 \\
\hline
人设 \& 自我呈现 & 角色扮演、个性特质、价值观和情绪状态。 \\
\hline
交互 \& 响应策略 & 对话导航、确定性/模糊表达、澄清策略、
征询行为和亲社会行为。 \\
\hline
结构 \& 推理 & 推理过程要求、论证结构、条件分支和范围约束。 \\
\hline
指令处理 & 条件性、多步指令链、冲突处理。 \\
\hline
\end{longtable}
}

大型仓库的构建时间。

主要的吞吐瓶颈是 LLM token 消耗，而非 CPU 或磁盘。该流水线运行在 10,000 个 CPU 核心上，具备每分钟 1200 万 tokens 的 LLM 容量，实现了 83\% 的提示词缓存命中率。在该配置下，流水线每分钟大约产出 20 个通过评分（grading）的环境。
